%% Formal semantics of the Encore VM (paper version).
%% Transcribed from crates/encore_vm/src/vm.rs and VM.md; keep in sync.
%% The full version, with the Scheme subset and the CPS IR, is
%% semantics-full.tex (standalone: encore-semantics.tex).

\newcommand{\kw}[1]{\mathsf{#1}}
\newcommand{\instr}[1]{\textsc{#1}}

\subsection{Formal Semantics of the VM}
\label{sec:semantics}

We give the semantics of the bytecode machine, following the standard treatment of compilation to a virtual machine~\cite{leroy2009mechanized}.
It is transcribed from the interpreter and not yet mechanized.
$h_j$ is the host function registered in extern slot $j$, which may read and allocate on the heap, and $\delta_{op}$ the meaning of a primitive opcode, a partial function undefined exactly where the VM traps (integer overflow, out-of-range byte).

\paragraph{Bytecode machine.}
Fig.~\ref{fig:vm} gives the abstract machine that \code{encore\_vm} implements.
It abstracts the 32-bit packed encoding of Fig.~\ref{fig:values} into tagged words $w$ and the arena into a partial map $H$ from addresses to objects, and leaves the code $c$ implicit: $c(pc)$ is the instruction at $pc$ and $pc^{+}$ the next one.
A state is $\langle pc, R, H \rangle$, with $R$ the 256-register file. As in the implementation, the globals live in the same memory as the heap: global $i$ is a reserved cell $g_i$ of $H$ holding a word.
There is no stack component: $\instr{encore}$ is the only transfer to a code address held in a value, and no instruction returns.
When allocation finds no free address, the collector replaces $H$ by $\pi(H|_{\mathit{reach}})$ and $R$ by $\pi(R)$, where $\mathit{reach}$ is the set of addresses reachable from $R$ and the global cells, and $\pi$ is an injective renaming of these addresses that fixes the $g_i$; the machine traps with \code{HeapOverflow} if allocation still fails.
The collector is thus invisible up to renaming, which is the property a proof of the mark-compact implementation must establish.

\begin{figure}[!tbp]
\centering\small
\[
\resizebox{\ifdim\width>\linewidth\linewidth\else\width\fi}{!}{$\displaystyle
\begin{array}{r@{\;}c@{\;}l}
w & ::= & \kw{int}(n) \mid \kw{func}(\ell) \mid \kw{clos}(a) \mid \kw{con}(t, a) \mid \kw{con}(t, \bullet) \mid \kw{bytes}(a) \\
H(a) & ::= & \langle \ell; w_0, \dots, w_{m-1} \rangle \ \text{(closure)} \mid \langle w_0, \dots, w_{m-1} \rangle \ \text{(fields)} \mid s \ \text{(byte string)} \mid w \ \text{(global cell)}
\end{array}
\qquad
\begin{array}{l}
\mathit{code}_H(\kw{func}(\ell)) = \ell \\
\mathit{code}_H(\kw{clos}(a)) = \ell \ \text{if}\ H(a) = \langle \ell; \bar{w} \rangle
\end{array}
$}
\]
\vspace{-0.3em}
\[
\resizebox{\ifdim\width>\linewidth\linewidth\else\width\fi}{!}{$\displaystyle
\begin{array}{@{}l@{\quad}l@{\quad}l@{}}
\toprule
c(pc) & \text{condition} & \text{next state} \\
\midrule
\instr{mov}\ d\ r & & \langle pc^+, R[d \mapsto R(r)], H \rangle \\
\instr{capture}\ d\ i & R(\kw{SELF}) = \kw{clos}(a),\ H(a) = \langle \ell; \bar{w} \rangle & \langle pc^+, R[d \mapsto w_i], H \rangle \\
\instr{global}\ d\ i & H(g_i) = w & \langle pc^+, R[d \mapsto w], H \rangle \\
\instr{function}\ d\ \ell & & \langle pc^+, R[d \mapsto \kw{func}(\ell)], H \rangle \\
\instr{closure}\ d\ \ell\ \bar{r} & a \notin \mathrm{dom}(H) & \langle pc^+, R[d \mapsto \kw{clos}(a)], H[a \mapsto \langle \ell; R(\bar{r}) \rangle] \rangle \\
\instr{pack}\ d\ t\ \bar{r} & \text{arity}(t) > 0,\ a \notin \mathrm{dom}(H) & \langle pc^+, R[d \mapsto \kw{con}(t, a)], H[a \mapsto \langle R(\bar{r}) \rangle] \rangle \\
\instr{pack}\ d\ t & \text{arity}(t) = 0 & \langle pc^+, R[d \mapsto \kw{con}(t, \bullet)], H \rangle \\
\instr{field}\ d\ r\ i & R(r) = \kw{con}(t, a),\ H(a) = \langle \bar{w} \rangle & \langle pc^+, R[d \mapsto w_i], H \rangle \\
\instr{unpack}\ d\ t\ r & R(r) = \kw{con}(t', a),\ H(a) = \langle \bar{w} \rangle & \langle pc^+, R[d{+}i \mapsto w_i]_{i < \text{arity}(t)}, H \rangle \\
\instr{match}\ r\ b\ \ell_0 \dots \ell_{m-1} & R(r) = \kw{con}(t, \_),\ 0 \le t - b < m & \langle \ell_{t-b}, R, H \rangle \\
\instr{encore}\ f\ k & \mathit{code}_H(R(f)) = \ell \in \mathrm{dom}(c) & \langle \ell, R[\kw{SELF} \mapsto R(f), \kw{CONT} \mapsto R(k)], H \rangle \\
\instr{extern}\ d\ r\ j & h_j(R(r), H) = (w, H') & \langle pc^+, R[d \mapsto w], H' \rangle \\
\instr{op}\ d\ \bar{r} & \delta_{op}(R(\bar{r}), H) = (w, H') & \langle pc^+, R[d \mapsto w], H' \rangle \\
\instr{fin}\ r & & \text{halt with } R(r) \\
\bottomrule
\end{array}
$}
\]
\caption{Bytecode machine (\code{vm.rs}). $\instr{branch}\ r\ b\ \ell_0\ \ell_1$ jumps to $\ell_0$ if the tag of $R(r)$ is $b$ and to $\ell_1$ otherwise; $\instr{global\_w}$ is $\instr{global}$ with a 16-bit index; $\instr{op}$ stands for the literal, integer, and byte-string opcodes, where $H'$ extends $H$ when the result is a new byte string. A state with no applicable rule is stuck. The interpreter traps on some stuck states (unknown tag, call outside the code, arithmetic overflow) but does not check value types: compiled code must never reach them.}
\label{fig:vm}
\end{figure}

\paragraph{Entry points.}
Code address $0$ holds the return stub $\instr{fin}\ \kw{A1}$.
The host calls a function value $f$ on arguments $\bar{w}$ by running the machine from $\langle \mathit{code}_H(f), R_0[\kw{SELF} \mapsto f, \kw{CONT} \mapsto \kw{func}(0), \kw{A}i \mapsto w_i], H \rangle$, with $R_0$ the cleared register file; the call returns $w$ when the machine halts with $w$.
Loading calls the entry point of each definition in order and stores the result in its global cell.

\paragraph{Correctness statement.}
Let $v \approx_H w$ relate first-order values to machine words: $n \approx_H \kw{int}(n)$; $s \approx_H \kw{bytes}(a)$ when $H(a) = s$; $C() \approx_H \kw{con}(t(C), \bullet)$; and $C(\bar{v}) \approx_H \kw{con}(t(C), a)$ when $H(a) = \langle \bar{w} \rangle$ and $v_i \approx_H w_i$ for all $i$, where $t(C)$ is the tag of constructor $C$.

\begin{conjecture}[Compiler correctness]
\label{conj:correct}
Let $P$ be an extracted program accepted by the frontend and $B$ the bytecode \encore{} produces for it.
If a definition $x$ of $P$ evaluates, under the call-by-value semantics of Scheme, to a first-order value $v$ (built from integers, byte strings, and constructors), then loading $B$ either initializes the global cell $g_x$ of $x$ to a word $w$ such that $v \approx_H w$ in the resulting heap $H$, or traps with \code{HeapOverflow}, \code{IntOverflow}, or \code{ByteRange}.
Likewise, if $x$ denotes a function and $x\ v_1 \dots v_m$ evaluates to a first-order $v$ for first-order arguments $v_i \approx_H w_i$, then calling $x$ on $\bar{w}$ returns a word $w$ with $v \approx_{H'} w$ in the final heap $H'$, or traps likewise.
\end{conjecture}

\noindent
Since the machine is deterministic, this forward simulation also excludes any other observable behavior~\cite{leroy2009mechanized}.
The statement decomposes along the pipeline into three simulations, each with a precedent in a mechanized compiler:
(i)~the frontend and the CPS transformation, as CertiRocq does for its front-end passes such as closure conversion~\cite{paraskevopoulou2019closure};
(ii)~each optimization pass, as in Pilsner~\cite{neis2015pilsner};
(iii)~register allocation and emission against the machine of Fig.~\ref{fig:vm}, as CertiRocq does for C~\cite{savarybelanger2019cpstoc}.
Two links remain outside the statement: Rocq's extraction, for which a verified alternative exists for OCaml~\cite{forster2024verified}, and the agreement between Fig.~\ref{fig:vm} and the Rust interpreter, discussed in Sect.~\ref{sec:discussion}.

\FloatBarrier
